% =========================
% FICHIER PRINCIPAL (main.tex)
% Classe: report (chapitres)
% Rôle: définir le préambule + assembler les chapitres
% =========================

\documentclass[12pt,a4paper]{report}
% \documentclass définit le type de document.
% report = structure avec \chapter, \section, etc.
% Les options (12pt, a4paper) règlent la taille de police et le format papier.

% ---------- Préambule commun ----------
\input{Preamble}
% \input{...} insère le contenu d’un fichier à cet endroit (comme un copier-coller).
% On l’utilise souvent pour mettre les packages et macros dans un fichier séparé.

% ---------- Compilation partielle (optionnel) ----------
% \includeonly{chapters/02-method,chapters/03-results}
% \includeonly permet de compiler uniquement certains chapitres
% (ceux appelés via \include) pour accélérer la compilation.
% Important: ça n’a d’effet que sur \include, pas sur \input.

\begin{document}
% Début du contenu visible du document.
% Tout ce qui est avant = préambule (configuration), tout ce qui est après = contenu.

% ---------- Pages de début (exemples) ----------
\title{Titre du rapport}
% Définit le titre (utilisé par \maketitle).

\author{Prénom NOM}
% Définit l’auteur (utilisé par \maketitle).

\date{\today}
% Définit la date (ici: date du jour). Peut être vide: \date{}.

\maketitle
% Génère la page de titre à partir de \title, \author, \date.

\tableofcontents
% Génère la table des matières (sommaire).

% ---------- Chapitres ----------
\include{Chapters/01_Intro/01_Intro}
% \include insère un “gros bloc” (souvent un chapitre).
% Particularité: force des sauts de page avant/après (pratique pour des chapitres).

\chapter{Marketing}
\section{Analyse Pestel}
L'analyse marketing a été structurée en deux grandes parties. D'une part, nouvs avons collecté un maximum de données concernant les besoins et attentes des patient en commencant par une analyse Pestel détaillée. Cette analyse nous a permis d'établir le contexte de marché. Les grands éléments à retenir pour cette analyse sontes les suivantes
 
\begin{itemize}
    \item Politique: Le développement du liner dépendra des subsides accordées par la région Wallonne. Ces subsides seront accordés suite à une compétition de financement publique auquel différents instituts (universités/hautes écoles) présenteront leurs projets. Le dossier consituté pour la région devra donc être particulièrement sourcé et solide
    \item Economique: Le marché mondial des prothèses de membres inférieurs est en forte croissance et devrait atteindre 2,6 milliards USD en 2030 contre 1,7 milliards USD en 2024. En plus d'un marché croissant, le système de remboursement INAMI offre un cadre favorable pour un tel marché avec des remboursements allant jusqu'à 2696,83€ pour un liner adaptatif. 
    \item Socioculturel: Les patients amputés font fassent à des douleurs et gènes au cours de la journée. La dilatation et retraction du moignon est incontrolable et oblige les patients à faire des ajustements pour limiter les gènes. Un dispositif comme un liner adaptatif semble être une solution novatrice jamais expérimentée auparavant.
    \item Technologique: Le futur liner utilisera des technologies de pointeavec capteurs de pression piézorésistifs ou capacitifs intégrés avec électrodes textiles soiples et lavables. Ces solutions sont développés par des entreprises et instituts européenne comme AG-tech (France) et les universités de Nantes et de Louvain (Louvain Bionics). Les défos technologique sont cependant majeurs avec une gestion de batterie qui devra être cruciale et une miniaturisation des composants qui ne doivent créer aucune gène chez le patient.
    \item Environnemental: Les ressources utilisées devront être le plus local possible car la région Wallonne y fait particulièrement attention. Malgré cela, certains composants ne pourront pas être exclusivement produits en Belgique. De plus, les contraintes de stockages des produits biomédicaux doivent respecter des températures et taux d'humidités précis.
    \item Légal: Le liner nouvellement développée devra respecter des normes et caractéristiques très précises afin d'être classifié comme un produit personnalisé. La position des capteurs de pression doit être à tout prix personnalisée pour que le produit puisse être vendu pour ~2700€. Les données collectées chez les patients doivent évidemment rester confidentielles en vertue de la protection des données personnelles (RGPD).
\end{itemize}
 
Cette analyse Pestel recensent le cadre dans lequel évoluera l'entreprise nouvellement crée et sera utilisée par la suite pour identifier les opportunités et menaces de la matrice SWOT.
\section{Exploration du marché (FIELD)}
 
L'exploration du marché a été réalisée en deux grandes parties, la première étant la collecte de donnée prilaire qualitatives. Pour réaliser cette collecte, un sondage de 33 questions a été réalisé et partager avec Mr.Fischer qui a formulé différente recommandations pour préciser les questions et assurer que le sondage collecte toutes les données nécéssaires auprès des patients. Une fois le sondage finalisé, celui-ci a été transmis une nouvelle fois et est actuellemnt diffusé auprès des clients de Tecfit qui donne leurs avis, expriment leurs espoirs/craintes par rapport à un liner automatique et soumettent leurs informations personnelles pour une interview en présentiel par la suite.  
 
\section{Exploration du marché (DESK)}
Pour les données secondaires, un grand nombre de sources ont été explorées afin de déterminer la taille du marché des liners en Belgique et en Europe de l'ouest. La première tâche a été de connaître l'évolution des amputation majeurs des membres inférieurs en Belgique. Les amputations majeurs sont celles réalisées au dessus de la cheville (tibia, genou, cuisse) et constitue le marché des patients susceptibles d'utiliser un liner. Deux sources ont été consultées, et fournissent l'évolution des amputations entre 2009 à 2013 et 2009 à 2018 pour le deuxième article. Une droite de régression a ensuite été tracée pour estimer les amputations de cette année et voici les résultats obtenus:
 
\begin{figure}[H]
    \centering
    \includegraphics[width=1\linewidth]{Fig_Marketing/Projection_2009_2018.png}
    \caption{Projection du nombre d'amputations}    \label{fig:Projetction_2009_2018}
\end{figure}


Nous constatons un taux d'amputations qui est en baisse à raison de 25 amputations en moins par an. Naturellement, ces chiffres se base sur la décénie passée et un plateau sera atteint car il n'est pas possible de ne plus avoir recourt à l'amputation, et ce malgré l'avancée de la médecine. Nous estimons donc que le nombre d'amputations pour 2025 se situe entre 950 et 1000. Ceci nous informe sur la tendance du marché, mais pas sur le nombre total de patients susceptiblent d'avoir recourt à un liner.
Il est très dur de connaître la taille du marché au patient près, un chiffre datant de 2022 nous informe qu'il y a +-7700 amputées en Belgiques (toutes amputations confondues). Ce chiffre se base sur une extrapolation du nombre d'amputés au Pays-Bas recalculé pour la population belge.
Le Pays-Bas étant un pays limitrophe de la Belgique avec un indice de développement humain similaire, ce chiffre de 7700 constitue une estimation solide. Sur base des différents articles collectés (Allemagne, Belgique, France, Pays-bas,...), nous constatons que les amputation des membre inférieurs représentent une minorité de cas et compte pour 5\% du nombre d'amputations total.
7300 est donc le nombre d'amputés de la jambe en Belgique. Sur base de ces même documents, le nombre d'amputations majeurs de la jambe par rapport au nombre d'amputations mineurs varie de 28\% jusqu'à 40 \% en fonction des sondages et des années. 33\% est donc un pourcentage moyen. Le nombre de patient avec amputations majeurs de la jambe en Belgique se situe entre 2400 et 2500 selon nos calculs (tous types confondus: diabétiques type 1, diabétiques type 2, accidents de travail).
Cela représente donc le marché maximal pouvant être atteint en Belgique. \\
 
Ce chiffre de 7700 bien que cité dans beaucoup de sources n'est pas justifié et aucun article ne semble expliquer comment il a été obtenu. Nous avons effectué une deuxième estimation du marché en se basant sur le taux de décès par an dans la population touchée par une amputation majeure. Plusieurs sources indiqiquent un taux de décès à 5 ans de 69\% pour les personnes diabétiques et 60\% pour les personnes non diébatiques. Cela nous donne un taux de survie à 5 ans de 35\%. 
\\
Taux de survie à 5 ans:
\begin{equation}
S_5 = 1 - 0,65 = 0,35
\end{equation}
Taux de mortalité par an:
\begin{equation}
S_5 = (1 - M)^5 => M = 1 - 0,35^\frac{1}{5} = 0,19
\end{equation}
 
Rapporté sur une année, le taux de décès est de 19\% de la population avec amputation majeure. Avec les nouveaux amputées précédemment calculer avec la régression précédente, voici les chiffres projetés pour 2025:
 
\begin{figure}[H]
    \centering
    \includegraphics[width=1\linewidth]{Fig_Marketing/Chiffres_Decenies.png}
    \caption{Projection du nombre d'amputés sur la décénie}    \label{fig:Chiffres_Decenies}
\end{figure}
 
L'ordre de grandeur par rapport au calcul prédédent est similaire, mais la valeur du marché total passe du simple au double. Pour conserver une approche conservatrice dans l'estimation du marché, la première valeur a été retenue. Nous etimerons donc que le nombre de personnes avec une amputation majeures en Belgique est de 2500.
 
\section{Besoin en flux de trésorerie}
Suite à une estimation des coûts fixes et variables sur 5 ans effectuée dans la partie logistique, une première courbe des besoins en flux de trésorerie a été réalisée. Cette courbe a pour but de mettre en évidence le besoin maximal en trésorerie et quand ce besoin maximal est atteint. En considérant une évolution linéaire (les patients changent de liner tous les ans) voici les résultats obtenus:
 
\begin{figure}[H]
    \centering
    \includegraphics[width=0.8\linewidth]{Fig_Marketing/Courbe_flux_tresorerie.png}
    \caption{Projection du nombre d'amputés sur la décénie}    \label{fig:Courbe_flux_tresorerie}
\end{figure}
 
\begin{itemize}
    \item Le besoin maximal est de 600000€ a lieu au bout de 25 mois 
    \item La rentabilité est atteinte au bout de 50 semaines
    \item Le nombre de pacients est croissant à raison de 3 nouveaux liners fabriqués tous les deux mois
\end{itemize}




\chapter{Finance}

À partir des bilans de Tecfit pour les exercices 2022, 2023 et 2024, une analyse financière des ratios clés a été réalisée afin d’évaluer la stabilité financière de l’entreprise et sa capacité à mener à bien un projet de recherche d’envergure.

L’étude s’appuie sur dix ratios financiers issus de la liste de la Banque nationale de Belgique (BNB). Pour chacun d’entre eux, nous présentons brièvement son intérêt, son évolution au fil des années, ainsi qu’une comparaison avec un acteur du même secteur.

L’entreprise de référence retenue est Aqtor, spécialisée dans la fabrication de prothèses et d’orthèses. Active depuis 2018 et intégrée au groupe Eqwal Ability Group, Aqtor opère également en Belgique, dans la région de Gand. Ce choix se justifie par la proximité d’activité entre les deux sociétés et la pertinence de la comparaison sectorielle.

Les données 2023 d’Aqtor ont été utilisées, car elles offrent une base d’informations complète et récente, nécessaire pour assurer la fiabilité des comparaisons.

\section{Marge Nette sur vente}
Le premier ratio analysé est la marge nette sur ventes, également appelée taux de valeur ajoutée dans la nomenclature de la Banque nationale de Belgique (BNB). Ce ratio mesure la proportion du bénéfice d’exploitation par rapport au chiffre d’affaires, autrement dit, la part du chiffre d’affaires qui demeure après déduction de l’ensemble des charges opérationnelles. On a préféré utiliser la marge nette sur vente plutôt que la marge brute sur vente car les amortissements comptables ne tiennent pas compte de la totalité du matériel nécessaire à l’activité. Par exemple, Tecfit dispose d’une pompe à vide achetée d’occasion qui a une vingtaine d’année et qui peut tenir encore longtemps car peu utilisée. Son amortissement comptable reflète donc mal son utilisation réelle.  
L’évolution du ratio sur les trois dernières années est la suivante :

\begin{table}[h]
\centering
\begin{tabular}{c|c|c}
\textbf{2022} & \textbf{2023} & \textbf{2024} \\
\hline
0,29 & 0,33 & 0,28 \\
\end{tabular}
\end{table}

Ces chiffres indiquent qu’en 2024, pour 1euro de ventes, 24 centimes ont été transformés en bénéfice d’exploitation. Malgré une légère baisse par rapport à 2023, ce niveau reste globalement satisfaisant et nettement supérieur à celui du concurrent direct, dont la marge nette s’élève à 0,09.
Cette différence traduit une rentabilité opérationnelle solide et une bonne maîtrise des coûts, suggérant que l’entreprise parvient à exploiter efficacement ses ressources pour générer du profit. La hausse observée en 2023 s’explique par un bond du bénéfice d’exploitation de 57\%, plus important que celui du chiffre d’affaires (+40\%), liée à une optimisation des charges (diminution du compte 60/61 mais augmentation du CA entre 2022 et 2023).
En revanche, la baisse de 2024 mérite attention: elle résulte d’une augmentation du chiffre d’affaires (+15\%) mais une légère diminution du bénéfice d’exploitation (-1.2\%).

\subsection{Taux de valeur ajoutée}
Ce ratio mesure la proportion entre la marge brute d’exploitation et la somme de la marge brute d’exploitation, des achats de matières premières et des services externes. Il reflète la capacité de l’entreprise à créer de la valeur à partir des ressources consommées pour sa production. Un ratio plus élevé traduit une utilisation plus efficace des charges externes et des matières premières.
L’évolution du ratio pour Tecfit est la suivante :

\begin{table}[h]
\centering
\begin{tabular}{c|c|c}
\textbf{2022} & \textbf{2023} & \textbf{2024} \\
\hline
0,27 & 0,44 & 0,45 \\
\end{tabular}
\end{table}
On observe une progression significative entre 2022 et 2023, avec un taux de valeur ajoutée passant de 0,27 à 0,44. Cette hausse s’explique principalement par une augmentation du chiffre d’affaires (de € 792 505 à € 1 108 783) combinée à une baisse des charges liées aux achats et services externes (compte 60/61 : de € 733 451,64 à € 661 167). Par conséquent, la marge brut d’exploitation a nettement progressé (de € 265,061.00 à € 511,992.00).
En 2024, le ratio se maintient à un niveau stable (0,45), traduisant une bonne efficacité opérationnelle constante. On note cependant une hausse des charges d’exploitation (compte 60/61 : + € 68 644 par rapport à 2023). Le maintien du ratio vient d’une augmentation identique de la marge d’exploitation et du compte 60/61, ils ont tous les 2 augmenté d’approximativement de 15\%. 
Par comparaison, Aqtor affichait en 2023 un ratio de 0,47, approximativement identique à celui de Tecfit. Cela confirme que Tecfit gère efficacement ses matières premières et services externes, dégageant une marge proportionnellement identique à celle de plus gros concurrent que lui.

\subsection{Valeur ajoutée par personne occupée}
Ce ratio mesure le bénéfice d’exploitation généré par chaque employé. Il permet d’évaluer la productivité du personnel et d’identifier si les effectifs sont dimensionnés de manière optimale. Cet indicateur est particulièrement pertinent pour Tecfit, dont l’activité repose sur une production artisanale à forte intensité de main-d’œuvre, sans recours à l’automatisation. La création de valeur dépend donc directement de la performance et de l’efficacité de l’équipe.
Évolution du ratio
2022: 155 918,06€  |2023: 124 876,83€|2024: 133 174,59€
\begin{table}[h]
\centering
\begin{tabular}{c|c|c}
\textbf{2022} & \textbf{2023} & \textbf{2024} \\
\hline
 155 918,06€ & 124 876,83€ & 133 174,59€ \\
\end{tabular}
\end{table}
En 2022, le ratio atteint un niveau particulièrement élevé, principalement en raison d’un nombre réduit d’employés équivalents temps plein (1,7ETP). En 2023, ce nombre passe à 4,1ETP, soit une hausse de 140\%, alors que la marge d’exploitation a progressé de 93\%. Cette évolution explique la baisse du bénéfice par employé, sans pour autant traduire une baisse d’efficacité, on remarque au ratio 2 que les couts du compte 60/61 sont relativement constant alors que le chiffre d’affaires augmente, cette meilleure utilisation des ressources matériel se fait au prix d’un peu plus de travail pour la main d’œuvre. 
La légère augmentation observée en 2024 (+6.6\%) s’explique par une meilleure intégration des nombreux nouveaux employés, encore en phase de formation ou d’adaptation en 2023. 
Comparaison sectorielle
En comparaison, Aqtor présentait en 2023 une marge brute d’exploitation par employé de 173 562€, soit un peu moins de 40\% de plus que les employés de Tecfit. Cela met en évidence qu’il y a une possibilité d’amélioration pour Tecfit, bien qu’il ne puisse probablement pas égaler les performances d’Aqtor qui bénéficie d’une meilleure spécialisation de ses employés au vu de la taille de la structure. 

\include{Chapters/03_Finance/03_Finance/analyse financière des ratios}
À partir des bilans de Tecfit pour les exercices 2022, 2023 et 2024, une analyse financière des ratios clés a été réalisée afin d’évaluer la stabilité financière de l’entreprise et sa capacité à mener à bien un projet de recherche d’envergure.

L’étude s’appuie sur dix ratios financiers issus de la liste de la Banque nationale de Belgique (BNB). Pour chacun d’entre eux, nous présentons brièvement son intérêt, son évolution au fil des années, ainsi qu’une comparaison avec un acteur du même secteur.

L’entreprise de référence retenue est Aqtor, spécialisée dans la fabrication de prothèses et d’orthèses. Active depuis 2018 et intégrée au groupe Eqwal Ability Group, Aqtor opère également en Belgique, dans la région de Gand. Ce choix se justifie par la proximité d’activité entre les deux sociétés et la pertinence de la comparaison sectorielle.

Les données 2023 d’Aqtor ont été utilisées, car elles offrent une base d’informations complète et récente, nécessaire pour assurer la fiabilité des comparaisons.

\include{subchapter/Marge Nette sur vente}
Le premier ratio analysé est la marge nette sur ventes, également appelée taux de valeur ajoutée dans la nomenclature de la Banque nationale de Belgique (BNB). Ce ratio mesure la proportion du bénéfice d’exploitation par rapport au chiffre d’affaires, autrement dit, la part du chiffre d’affaires qui demeure après déduction de l’ensemble des charges opérationnelles. On a préféré utiliser la marge nette sur vente plutôt que la marge brute sur vente car les amortissements comptables ne tiennent pas compte de la totalité du matériel nécessaire à l’activité. Par exemple, Tecfit dispose d’une pompe à vide achetée d’occasion qui a une vingtaine d’année et qui peut tenir encore longtemps car peu utilisée. Son amortissement comptable reflète donc mal son utilisation réelle.  
L’évolution du ratio sur les trois dernières années est la suivante :

\begin{table}[h]
\centering
\begin{tabular}{c|c|c}
\textbf{2022} & \textbf{2023} & \textbf{2024} \\
\hline
0,29 & 0,33 & 0,28 \\
\end{tabular}
\end{table}

Ces chiffres indiquent qu’en 2024, pour 1 euro de ventes, 24 centimes ont été transformés en bénéfice d’exploitation. Malgré une légère baisse par rapport à 2023, ce niveau reste globalement satisfaisant et nettement supérieur à celui du concurrent direct, dont la marge nette s’élève à 0,09.
Cette différence traduit une rentabilité opérationnelle solide et une bonne maîtrise des coûts, suggérant que l’entreprise parvient à exploiter efficacement ses ressources pour générer du profit. La hausse observée en 2023 s’explique par un bond du bénéfice d’exploitation de 57\%, plus important que celui du chiffre d’affaires (+40\%), liée à une optimisation des charges (diminution du compte 60/61 mais augmentation du CA entre 2022 et 2023).
En revanche, la baisse de 2024 mérite attention : elle résulte d’une augmentation du chiffre d’affaires (+15\%) mais une légère diminution du bénéfice d’exploitation (-1.2\%).

\include{subchapter/Taux de valeur ajoutée}
Ce ratio mesure la proportion entre la marge brute d’exploitation et la somme de la marge brute d’exploitation, des achats de matières premières et des services externes. Il reflète la capacité de l’entreprise à créer de la valeur à partir des ressources consommées pour sa production. Un ratio plus élevé traduit une utilisation plus efficace des charges externes et des matières premières.
L’évolution du ratio pour Tecfit est la suivante :

\begin{table}[h]
\centering
\begin{tabular}{c|c|c}
\textbf{2022} & \textbf{2023} & \textbf{2024} \\
\hline
0,27 & 0,44 & 0,45 \\
\end{tabular}
\end{table}
On observe une progression significative entre 2022 et 2023, avec un taux de valeur ajoutée passant de 0,27 à 0,44. Cette hausse s’explique principalement par une augmentation du chiffre d’affaires (de € 792 505 à € 1 108 783) combinée à une baisse des charges liées aux achats et services externes (compte 60/61 : de € 733 451,64 à € 661 167). Par conséquent, la marge brut d’exploitation a nettement progressé (de € 265,061.00 à € 511,992.00).
En 2024, le ratio se maintient à un niveau stable (0,45), traduisant une bonne efficacité opérationnelle constante. On note cependant une hausse des charges d’exploitation (compte 60/61 : + € 68 644 par rapport à 2023). Le maintien du ratio vient d’une augmentation identique de la marge d’exploitation et du compte 60/61, ils ont tous les 2 augmenté d’approximativement de 15\%. 
Par comparaison, Aqtor affichait en 2023 un ratio de 0,47, approximativement identique à celui de Tecfit. Cela confirme que Tecfit gère efficacement ses matières premières et services externes, dégageant une marge proportionnellement identique à celle de plus gros concurrent que lui.

\include{subchapter/Valeur ajoutée par personne occupée}
Ce ratio mesure le bénéfice d’exploitation généré par chaque employé. Il permet d’évaluer la productivité du personnel et d’identifier si les effectifs sont dimensionnés de manière optimale. Cet indicateur est particulièrement pertinent pour Tecfit, dont l’activité repose sur une production artisanale à forte intensité de main-d’œuvre, sans recours à l’automatisation. La création de valeur dépend donc directement de la performance et de l’efficacité de l’équipe.
Évolution du ratio
2022 : 155 918,06 €   |  2023 : 124 876,83 €  |  2024 : 133 174,59 €
\begin{table}[h]
\centering
\begin{tabular}{c|c|c}
\textbf{2022} & \textbf{2023} & \textbf{2024} \\
\hline
 155 918,06 € & 124 876,83 € & 133 174,59 € \\
\end{tabular}
\end{table}
En 2022, le ratio atteint un niveau particulièrement élevé, principalement en raison d’un nombre réduit d’employés équivalents temps plein (1,7 ETP). En 2023, ce nombre passe à 4,1 ETP, soit une hausse de 140\%, alors que la marge d’exploitation a progressé de 93\%. Cette évolution explique la baisse du bénéfice par employé, sans pour autant traduire une baisse d’efficacité, on remarque au ratio 2 que les couts du compte 60/61 sont relativement constant alors que le chiffre d’affaires augmente, cette meilleure utilisation des ressources matériel se fait au prix d’un peu plus de travail pour la main d’œuvre. 
La légère augmentation observée en 2024 (+ 6.6 \%) s’explique par une meilleure intégration des nombreux nouveaux employés, encore en phase de formation ou d’adaptation en 2023. 
Comparaison sectorielle
En comparaison, Aqtor présentait en 2023 une marge brute d’exploitation par employé de 173 562€, soit un peu moins de 40\% de plus que les employés de Tecfit. Cela met en évidence qu’il y a une possibilité d’amélioration pour Tecfit, bien qu’il ne puisse probablement pas égaler les performances d’Aqtor qui bénéficie d’une meilleure spécialisation de ses employés au vu de la taille de la structure. 

% \include{chapters/03-results}

% ---------- Annexes ----------
\appendix
% À partir d’ici, les chapitres deviennent des annexes (A, B, C... au lieu de 1, 2, 3).

% \include{appendices/app-a}

% ---------- Fin du document ----------
\end{document}
% Fin du document: rien ne doit apparaître après.

